\VerdierSemanticAnchor{front-editors-note}
\VerdierPhysicalPageAnchor{0008}
\chapter*{Note des éditeurs}
\addcontentsline{toc}{chapter}{Note des éditeurs}

\VerdierSemanticAnchor{front-editors-note-publication-history}
C'est lors de la cérémonie en hommage à J.-L. Verdier, organisée par
l'Université Paris VII, le 19 octobre 1989, que M. Artin, P. Cartier et
L. Illusie eurent l'idée de publier la thèse de J.-L. Verdier. Deux ans plus
tard je me joignais à eux et j'obtenais, grâce à l'appui de J. Le Potier et de
M. Broué, un financement de Paris VII (URA 212) et de l'ÉNS (DMI) pour une
saisie de ce texte en \TeX{}. Une fois la frappe terminée, le résultat fut
soigneusement comparé à l'original dactylographié de Verdier. Il restait encore
à faire une relecture mathématique car Verdier n'avait jamais effectué les
dernières corrections.

\VerdierSemanticAnchor{front-editors-note-chapter-i-modifications}
Cette lecture attentive a été faite, avec un extrême soin, par P. Cartier
(chapitre I), B. Keller (chapitre III), et, pour l'ensemble du texte, par
M. Zisman et moi-même. En dehors des nombreuses coquilles qui ne méritent pas
d'être énumérées, les principales modifications apportées sont les suivantes.
Dans le chapitre I, la proposition (3.1.5) affirmait une équivalence. Cette
équivalence a été remplacée par une implication, la réciproque (jamais utilisée
par la suite) étant fausse. Dans la proposition (3.4.1), il y avait une erreur
de signe, induisant des erreurs de signe dans le chapitre III ((1.3.4), (b) et
(3.2.7.3)). Ceci a conduit à adapter la définition de la “catégorie triangulée
opposée” (chap. II, 1.1.7) afin que reste vraie l'affirmation :
\emph{la catégorie dérivée de la catégorie opposée est la catégorie triangulée
opposée à la catégorie dérivée}.

\VerdierSemanticAnchor{front-editors-note-chapter-ii-modifications}
Dans le chapitre II, l'ordre des propositions du §1.2 a été modifié. En effet,
Verdier s'était rendu compte, après la rédaction, que l'assertion que la somme
directe de deux triangles distingués est un triangle distingué est conséquence
des autres axiomes des catégories triangulées. Il l'avait donc supprimée de la
liste des axiomes, sans effectuer les remaniements qui en résultaient. D'autre
part, les alinéas (2.3.6) et (2.3.7) qui n'étaient pas utilisés dans la suite,
et dont la vérification a échoué, ont été supprimés. La définition de la notion
d'objet spectral stationnaire (4.4.2) a été légèrement modifiée, la définition
originale étant trop faible pour l'établissement du théorème (4.4.3) qui
suivait. Il a été vérifié que cela était sans conséquence pour la suite.

\VerdierPhysicalPageAnchor{0009}
\VerdierSemanticAnchor{front-editors-note-apparatus}
L'index de notations et l'index terminologique ont été complétés. En revanche,
la bibliographie a été laissée inchangée (sauf pour les SGA dont les références
définitives dans les “Lecture Notes” ont été indiquées). En particulier,
l'absence de référence numéro 13 n'a pas été comblée.

\VerdierSemanticAnchor{front-editors-note-acknowledgements}
Je voudrais remercier vivement M. Broué, A. Bruguières, P. Cartier, L. Gruson,
L. Illusie, B. Keller, J. Le Potier, M. Zisman, ainsi que l'Université Paris
VII, l'École Normale Supérieure, et l'Institut des Hautes Études Scientifiques,
qui ont aidé à divers titres à la réalisation de ce projet. Je remercie
également C. Defosse et C. Gourgues pour avoir patiemment saisi cet ouvrage en
\textsc{AMS}-\TeX{}.

\VerdierSemanticAnchor{front-editors-note-memorial}
J'aurais aimé pouvoir étendre ces remerciements à L. Doustaing prématurément
disparu en décembre 1994.

\VerdierSemanticAnchor{front-editors-note-signature}
\begin{flushright}
Paris, le 14 juin 1996\\
Georges Maltsiniotis
\end{flushright}

